\section{Small orbits give supports at every rank}\label{sec:support}\label{sup:support}

We now begin the lower-bound argument. The group $H$ of automorphisms from
Section~\ref{str:structures} has an abelian quotient of face translations
and a central subgroup $K$ of edge circulations. Commutators connect these
two parts. We use that connection to show that an object with a small
$H$-orbit is fixed by the pointwise stabilizer in $K$ of a short tuple of
atoms. The bound does not depend on the object's rank. We give the full
support argument from~\cite[Section~3, Proposition~8]{cptCompanion2026}.

Fix $L\ge2$ and $b\in\F^{V_L}$. Write $V=V_L$, $E=E_L$, $F=F_L$,
and $N=|A_b|$. Recall the incidence map
$\partial:\F^E\to\F^V$, the oriented face boundaries $c^f$, and
\[
 K=\ker\partial,\qquad C_F=(\F^2)^F,
 \qquad H=C_F\times K
\]
with the group law of Section~\ref{str:structures}. We identify $K$ with
the central subgroup $\{0\}\times K$ of $H$. Its action fixes face
coordinates and adds $k_e$ to the scalar coordinate of every edge atom
on $e$. As usual, an atom permutation acts on finite sets recursively by
$h\cdot x=\{h\cdot y:y\in x\}$.

\begin{definition}\label{sup:definition}
For an atom tuple $\alpha=(\alpha_1,\ldots,\alpha_r)$, let
\[
 K_\alpha=\{k\in K:k\cdot\alpha_j=\alpha_j
                         \text{ for every }1\le j\le r\}.
\]
The tuple $\alpha$ is a \emph{$K$-support} of an atom or hereditarily
finite set $x$ over $A_b$ if every element of $K_\alpha$ fixes $x$.
The empty tuple is permitted, and $K_{()}=K$.
\end{definition}

The proof has two steps. First, an edge functional whose curl is nonzero
on few faces can be represented on few edges. Second, the small index of
an object's stabilizer forces precisely this curl condition for the
functionals that detect its central stabilizer.

\subsection{Sparse representatives of edge functionals}

An \emph{edge cochain} is a vector $\ell\in\F^E$. Define its curl
$D\ell\in\F^F$ by
\[
 (D\ell)(f)=\sum_{e\in E}c^f_e\ell_e.
\]
For a vertex function $\varphi\in\F^V$, its gradient is
$\partial^{\mathsf T}\varphi$; on an edge directed from $v$ to $w$ it
has value $\varphi(w)-\varphi(v)$. Since $\partial c^f=0$, adding a
gradient changes neither the curl nor the functional
$k\mapsto\sum_e\ell_e k_e$ on $K$.

\begin{lemma}[Sparse cochains]\label{sup:sparse-cochain}
Let $\ell\in\F^E$, and let
$t=|\{f\in F:(D\ell)(f)\ne0\}|$. There is a vertex function
$\varphi\in\F^V$ such that
\[
 |\operatorname{supp}(\ell-\partial^{\mathsf T}\varphi)|\le2Lt,
 \qquad
 \operatorname{supp}(r)=\{e\in E:r_e\ne0\}.
\]
In particular, every cochain with zero curl is a gradient.
\end{lemma}

\begin{proof}
We first eliminate the cochain in one coordinate direction. We then
eliminate a second direction on a slice containing few nonzero curls,
and use the remaining curls to bound the support away from that slice.

Write vertices as $(i,j,k)$. The coordinate $\ell_1(i,j,k)$ belongs to
the edge from $(i,j,k)$ to $(i+1,j,k)$, and $\ell_2,\ell_3$ are defined
similarly. A direction-$a$ edge coordinate is used only when increasing
coordinate $a$ stays in the box. Put
\[
 \varphi_1(i,j,k)=\sum_{u=0}^{i-1}\ell_1(u,j,k),
 \qquad r^{(1)}=\ell-\partial^{\mathsf T}\varphi_1,
\]
where empty sums are zero. Then $r^{(1)}_1=0$. The face orientation
$+1,+a,-1,-a$ gives
\begin{align}
 (D\ell)_{12}(i,j,k)
   &=r^{(1)}_2(i+1,j,k)-r^{(1)}_2(i,j,k),\label{sup:curl12}\\
 (D\ell)_{13}(i,j,k)
   &=r^{(1)}_3(i+1,j,k)-r^{(1)}_3(i,j,k).\label{sup:curl13}
\end{align}
Here $ab$ indicates the face in directions $a<b$ with lower corner
$(i,j,k)$, and the identity is asserted whenever that face exists.

Let $t_{ab}$ be the number of nonzero $ab$ curls, so
$t=t_{12}+t_{13}+t_{23}$. The $23$ faces are partitioned into the $L$
slices of fixed first coordinate. Choose a slice $i=i_0$ containing
$t_0\le t_{23}/L$ nonzero $23$ curls. Define
\[
 \varphi_2(i,j,k)=\sum_{v=0}^{j-1}r^{(1)}_2(i_0,v,k),
 \qquad r^{(2)}=r^{(1)}-\partial^{\mathsf T}\varphi_2.
\]
Because $\varphi_2$ is independent of $i$, we still have
$r^{(2)}_1=0$, and now $r^{(2)}_2(i_0,j,k)=0$. Finally put
\[
 \varphi_3(i,j,k)=\sum_{w=0}^{k-1}r^{(2)}_3(i_0,0,w),
 \qquad r=r^{(2)}-\partial^{\mathsf T}\varphi_3.
\]
The function $\varphi_3$ is independent of $i,j$, so
\[
 r_1=0,\qquad r_2(i_0,j,k)=0,\qquad r_3(i_0,0,k)=0.
\]
All three gradient subtractions preserve curl. On the selected slice
the $23$ curl identity therefore becomes
\[
 r_3(i_0,j+1,k)-r_3(i_0,j,k)=(D\ell)_{23}(i_0,j,k).
\]
For each fixed $k$, every nonzero value of $r_3(i_0,j,k)$ has a nonzero
curl earlier on its $j$-line. One nonzero curl can affect at most $L$
positions on that line. Thus at most $Lt_0$ direction-$3$ edge
coordinates on the slice are nonzero. Direction $2$ is zero there.

Copy these direction-$2$ and direction-$3$ slice values to every first
coordinate $i$, obtaining a reference cochain supported on at most
$L^2t_0$ edges, with direction $1$ zero. Identities
\eqref{sup:curl12} and~\eqref{sup:curl13} remain valid for $r$.
Consequently
\[
 r_2(i,j,k)-r_2(i_0,j,k)
\]
is the signed sum of the $12$ curls between $i_0$ and $i$ on that line.
A nonzero $12$ curl occurs in at most $L$ such sums, so at most
$Lt_{12}$ direction-$2$ coordinates differ from the reference cochain.
The same reasoning bounds the differing direction-$3$ coordinates by
$Lt_{13}$. A coordinate in the support of $r$ must either belong to
the reference support or differ from its reference value. Hence
\[
 |\operatorname{supp}(r)|
 \le L^2t_0+L(t_{12}+t_{13})\le Lt\le2Lt.
\]
Taking $\varphi=\varphi_1+\varphi_2+\varphi_3$ proves the claim. The
same construction gives $r=0$ when $t=0$.
\end{proof}

\subsection{From a small orbit to a support}

\begin{proposition}[Small-orbit support bound]\label{sup:small-orbit}
Let $L\ge2$, $b\in\F^{V_L}$, and $N=|A_b|$. Let $q\ge0$ and let $x$
be an atom or a hereditarily finite set over $A_b$. If
$|H\cdot x|\le N^q$, then $x$ has a $K$-support consisting of at most
\[
 2L(q\log_3N)^2
\]
edge atoms. There is no bound on the finite rank of $x$ in this statement.
\end{proposition}

\begin{proof}
Let $S=\{h\in H:h\cdot x=x\}$ be the stabilizer of $x$. Let $U$ be
its image under the projection $H\to C_F$, and put $K_0=S\cap K$.
Both $U$ and $K_0$ are subgroups of additive groups over the prime field
$\F$, and hence are vector subspaces. Define
\[
 a=\operatorname{codim}_{C_F}U,\qquad
 d=\operatorname{codim}_{K}K_0,\qquad
 h=\log_3[H:S].
\]
The restriction of the projection to $S$ has image $U$ and kernel $K_0$.
Every fibre has size $|K_0|$, so $|S|=|U|\,|K_0|$. Since
$|H|=|C_F|\,|K|$, orbit--stabilizer gives
\begin{equation}\label{sup:index}
 a+d=h=\log_3|H\cdot x|\le q\log_3N.
\end{equation}

Consider the annihilator
\[
 W=\{\lambda\in K^*: \lambda(k)=0\text{ for all }k\in K_0\},
 \qquad \dim W=d,
\]
where $K^*$ is the linear dual of $K$. Each $\lambda\in K^*$ extends
to a linear functional on $\F^E$: extend a basis of $K$ to one of
$\F^E$ and assign arbitrary values on the additional basis vectors.
The standard coordinate pairing thus gives an edge cochain $\ell$ with
\[
 \lambda(k)=\sum_{e\in E}\ell_e k_e\quad(k\in K).
\]
Two representatives differ by an element of
$K^\perp=\operatorname{im}\partial^{\mathsf T}$. Indeed the latter
space is contained in $K^\perp$, and both have dimension
$\operatorname{rank}\partial$. In particular, their curls agree.

Fix $\lambda\in W$ and a representative $\ell$. For any $u,u'\in U$,
choose lifts in $S$. Their commutator belongs to $K_0$; the commutator
formula~\eqref{str:commutator} therefore implies
\begin{equation}\label{sup:isotropic}
 0=\sum_{f\in F}(D\ell)(f)\,\omega(u_f,u'_f).
\end{equation}
The right-hand side defines an alternating bilinear form $\beta_\ell$
on $C_F$. If $t$ of the coefficients $(D\ell)(f)$ are nonzero, its
rank is $2t$: each nonzero coefficient multiplies a nondegenerate
two-dimensional alternating block, and the face coordinates are
independent summands of $C_F$. Relations among the boundary vectors
$c^f$ do not identify these summands.

We use the standard dimension bound for isotropic subspaces of an
alternating form, after quotienting by its radical
\cite[Sections~3.1 and~3.3]{cameron2000classical}.
Equation~\eqref{sup:isotropic} says that $U$ is totally isotropic for
$\beta_\ell$, meaning that the form vanishes on $U\times U$. Put
$n_F=\dim C_F$, and let
$R=\{v\in C_F:\beta_\ell(v,w)=0\text{ for every }w\in C_F\}$
be the radical of this form. The quotient $C_F/R$ is nondegenerate of
dimension $2t$. The image $\overline U$ of $U$
is isotropic, so $\overline U\subseteq\overline U^{\perp}$ and
$\dim\overline U\le t$. The radical has dimension $n_F-2t$, whence
\[
 \dim U\le(n_F-2t)+t=n_F-t,
 \qquad t\le a.
\]
Lemma~\ref{sup:sparse-cochain} now replaces $\ell$ by a representative
of the same functional $\lambda$ supported on at most $2La$ edges.

We have obtained a sparse representative for every functional detecting
the quotient $K/K_0$. Choose such representatives for a basis of $W$,
and let $T\subseteq E$ be the union of their supports. Then
\begin{equation}\label{sup:edge-support-size}
 |T|\le2Lad\le2Lh^2\le2L(q\log_3N)^2.
\end{equation}
When $d=0$, use the empty basis and $T=\emptyset$. If $k\in K$
vanishes on $T$, every member of that basis vanishes on $k$. Thus all
of $W$ vanishes on $k$, and finite-dimensional duality gives $k\in K_0$.
For each $e\in T$, choose one edge atom on $e$, and list the chosen
atoms in a tuple $\alpha$. The central action in
\eqref{str:edge-action} fixes that atom exactly when $k_e=0$.
Consequently $K_\alpha\subseteq K_0\subseteq S$, so $\alpha$ supports
$x$ and satisfies~\eqref{sup:edge-support-size}. This argument concerns
the stabilizer of $x$ alone and makes no induction on its rank.
\end{proof}

\subsection{Hereditary support for finite invariant families}

A family $\mathcal X$ of atoms and hereditarily finite sets is
\emph{membership-transitive} if $y\in x\in\mathcal X$ implies
$y\in\mathcal X$. It then contains every object reached from a member
by any finite sequence of membership steps. We call an object
\emph{hereditarily $s$-supported} if it and every such recursive member
have a $K$-support of length at most $s$; these supports need not coincide.

\begin{corollary}\label{sup:hereditary}
Let $L\ge2$, $b\in\F^{V_L}$, and $N=|A_b|$. Suppose $\mathcal X$ is
a finite $H$-invariant, membership-transitive family of atoms and
hereditarily finite sets over $A_b$, with $|\mathcal X|\le N^q$ for
some $q\ge0$. Every member of $\mathcal X$ and every recursive member
has a $K$-support of length at most $2L(q\log_3N)^2$. In particular,
every member is hereditarily $s$-supported for
\begin{equation}\label{sup:length}
 s=\left\lceil2L(q\log_3N+1)^2\right\rceil+1.
\end{equation}
All the supports may be padded to length exactly $s$. For fixed $q$,
\[
 s=O_q\bigl(L(1+\log L)^2\bigr).
\]
\end{corollary}

\begin{proof}
$H$-invariance places the orbit of each $x\in\mathcal X$ inside
$\mathcal X$. Proposition~\ref{sup:small-orbit} therefore gives the
claimed support for every member. Membership transitivity gives the
same conclusion for all recursive members. Adding coordinates to an
atom tuple can only shrink its pointwise stabilizer, so it preserves
the support property. Since $L\ge2$, the base has an edge atom; append
copies of any one such atom to reach length $s$, including when the
original support is empty. Finally Lemma~\ref{str:size} gives
$N=\Theta(L^3)$, which yields the displayed asymptotic bound.
\end{proof}
